\def\level{Première}  % Classe à droite
\def\course{NSI}
\def\chaptername{Architecture Von Neumann} % Titre au centre
\def\docname{ARSE01}        % Identifiant à gauche

\pagestyle{cours_FP}
\makeheader{} % Affiche le bandeau


\begin{easybox}[Architecture Von Neumann]{}
	\begin{minipage}{0.7\linewidth}
		Le mathématicien John Von Neumann a élaboré en 1944 la première description d'un ordinateur dont le programme est stocké dans sa mémoire.

		Son architecture comporte 4 éléments:

		\begin{itemize}
			\item L' \textbf{Unité Arithmétique et Logique (UAL)} qui effectue les opérations de base.
			\item L'\textbf{Unité de Contrôle} qui est chargé de séquencer les opérations.
			\item La \textbf{mémoire} qui contient les données et les instructions
			\item Les \textbf{périphériques (entrées, sorties)} qui communiquent avec l'extérieur.
		\end{itemize}

	\end{minipage}\hfill
	\begin{minipage}{0.28\linewidth}
		\includegraphics[width=1\linewidth]{archi-vneumann.jpg}
	\end{minipage}
	
\end{easybox}


%\begin{easybox*}[commontheoremstyle, sidebyside, sidebyside align = top]{Opérations mathématiques de base en Python}{}

\begin{easybox}{La carte mère}{}
	C'est une carte de circuit imprimé qui permet de connecter tous les composants entre eux (processeur, périphériques I/O, mémoire etc.) On y retrouve:
	
	\begin{minipage}{0.49\linewidth}
		
	\begin{itemize}
		\item un processeur connecté sur le socket CPU.
		\item de la mémoire vive (RAM).
		\item connecteurs en façade (USB, HDMI, RJ45\dots{})
		\item une pile pour le BIOS
		\item connecteurs pour composants internes (disque dur, cartes supplémentaires\dots{}) 
	\end{itemize}

	\end{minipage}\hfill{}
	\begin{minipage}{0.49\linewidth}
	\includegraphics[width=1\linewidth]{carte_mere_interne.jpg}
	\end{minipage}
\end{easybox}

\begin{easybox}{Les connecteurs}{}

		\begin{minipage}{0.49\linewidth}
		
	\begin{itemize}
		\item des ports USB (Communication)
		\item des ports HDMI / VGA / DVI (Video)
		\item des ports Ethernet RJ45 (Réseau)
		\item des prises jack (audio)
	\end{itemize}
	
	\end{minipage}\hfill{}
	\begin{minipage}{0.49\linewidth}
	\includegraphics[width=0.9\linewidth]{peripherique.jpg}
	\end{minipage}


\end{easybox}


\begin{easybox}{Role du processeur}{}
	Le \textbf{microprocesseur} (Central Processing Unit: CPU) est chargé d'interpréter et d'exécuter les instructions d'un programme, de lire et sauvegarder les résultat dans une mémoire.

	Il est constitué de:
	\begin{itemize}
		\item une \textbf{horloge} qui synchronise toutes les opérations
		\item une \textbf{unité de commande} qui appelle les données et décode les instructions
		\item une \textbf{unité arithmétique et logique (UAL)} qui effectue les calculs et les comparaisons
		\item des \textbf{registres} qui stockent les opérandes et les résultats des calculs.
	\end{itemize}

\end{easybox}


\begin{easybox}{Les processeurs aujourd'hui}
\renewcommand{\arraystretch}{1.2}
		\begin{minipage}{0.49\linewidth}
		
	\begin{tabular}{|l|c|c|c|c|l|}
        \hline
        \textbf{Processeur}  & \textbf{Core} & \textbf{Fréq }  & \textbf{Usage} \\
        \hline
        AMD Ryzen9  & 16 & 4,3GHz  & Gaming +\\
        \hline
        Intel Core i9  & 24 & 3,2GHz  & Gaming \\
        \hline
        AMD Ryzen7   & 8 & 4,7GHz & Gaming \\
        \hline
        Intel Core i7  & 20 & 3,4GHz & Gaming \\
        \hline
        AMD Ryzen5   & 6 & 3,9GHz  & Bureautique, \\
        \hline
    \end{tabular}

	\end{minipage}\hfill{}
	\begin{minipage}{0.49\linewidth}
\includegraphics[width=0.95\linewidth]{arse01.png}
	\end{minipage}
\end{easybox}




\begin{easybox}{Mémoire vive et mémoire morte}{}

	\begin{itemize}
		\item La \textbf{mémoire vive} (RAM: ``Random Access Memory'') permet de stocker les informations pendant tout le temps de fonctionnement de l'ordinateur.
		
		Son contenu est par contre détruit dès lors que l'ordinateur est éteint ou redémarré, contrairement à une mémoire de masse telle que le disque dur, capable de garder les informations même lorsqu'il est hors tension. On dit que cette mémoire est \textbf{volatile}.
	
		Cette mémoire est bien plus rapide les périphériques de stockage de masse (Disque dur, clé USB \dots) 
	
		\item La \textbf{mémoire morte} (ROM: ``Read Only Memory'') permet de stocker des données nécessaires au démarrage de l'ordinateur.
		
		Son contenu ne s'efface pas lors de la mise hors tension du système. C'est une mémoire non volatile qui contient les éléments essentiels au démarrage de l'ordinateur (BIOS).
				
		\item La \textbf{mémoire de masse} (Disque dur~--~HDD) sert à conserver les données de manière permanente. 
	\end{itemize} 
	
\end{easybox}

\begin{easybox}{Vitesse des différentes mémoires}

	\begin{minipage}{0.49\linewidth}
		\begin{tabular}{|l|c|c|}
			\hline
			\textbf{Type de mémoire} & \textbf{Débit moyen}  \\
			\hline
			RAM (DDR4/DDR5) & 25--80 Go/s\\
			\hline
			SSD NVMe (PCIe 4.0/5.0) & 3--14 Go/s \\
			\hline
			SSD SATA & 500--600 Mo/s\\
			\hline
			HDD SATA & 80--200 Mo/s \\
			\hline
			ROM du BIOS (EEPROM) & 50--100 ns \\
			\hline
		\end{tabular}

	\end{minipage}\hfill{}
	\begin{minipage}{0.49\linewidth}
	\includegraphics[width = 1\linewidth]{memoire.jpg}
	\end{minipage}		
\end{easybox}

\begin{exercice_cours}{Reconnaitre les différents types de mémoires}{}

\begin{itemize}
	\item Associez chaque image à son type de mémoire (RAM, SSD NVMe, SSD SATA, HDD SATA, ROM BIOS).
\end{itemize}

\includegraphics[width=0.3\linewidth]{ram4.jpg}
\includegraphics[width=0.3\linewidth]{nvme.jpg}
\includegraphics[width=0.3\linewidth]{rom_bios.jpg}

\includegraphics[width=0.28\linewidth]{sata_ssd.jpg}
\includegraphics[width=0.28\linewidth]{sata_hdd.jpg}
\includegraphics[width=0.28\linewidth]{proc.jpg}
\end{exercice_cours}

